\section{Load-Disconnection Energy Control}
\label{sec:disc-open}

After load disconnection, stored resonant energy and any additional packets
can continue charging the secondary capacitance. For packet energy
\(\Delta E\) delivered to capacitance \(C_o\), the idealized voltage rise is
\begin{equation}
V_{o,\mathrm{after}}=\sqrt{V_{o,\mathrm{before}}^2+\frac{2\Delta E}{C_o}}.
\label{eq:disc-burst-voltage-rise}
\end{equation}
Packet energy can therefore produce a substantial voltage step even at
low average duty. Capacitor discharge and the response of the USB-C PD stage also affect
the voltage rise.

Startup precharge, a no-load duty cap, a voltage window, and soft
overvoltage holdoff were implemented. During load recovery, the no-load
restriction is bypassed below the recovery threshold to restore transfer.

The captures demonstrate no-load and load-removal operation. Quantifying
the protection margin requires synchronized peak secondary voltage, peak
VBUS, active thresholds, firmware state, and recovery timing.

An independent secondary shutdown or discharge path should also be added.
The existing regulation and protection share MCU execution, sensing,
calibration, and enable output, leaving common failure paths in the
regulating firmware.
